\documentclass[12pt,a4paper]{article}

%% PACHETE MATEMATICE
\usepackage{amsmath,amssymb,amsfonts}
\usepackage{amsthm}
\usepackage{mathpazo}
\usepackage{eulervm}
\usepackage{enumerate}
\usepackage{color}
\usepackage{graphicx}
\usepackage[all]{xy}
\usepackage{xcolor}
\usepackage{framed}
\usepackage[unicode]{hyperref}
\setcounter{MaxMatrixCols}{20}
\usepackage{multicol}
% \usepackage[romanian]{babel}


%% ASPECT
\linespread{1.05}
\usepackage[T1]{fontenc}
\usepackage[utf8x]{inputenc}
\usepackage[margin=2cm]{geometry}
% \usepackage[color]{showkeys}
\usepackage{anyfontsize}
\usepackage{titlesec}
\titleformat*{\section}{\large\bfseries}

\usepackage{fancyhdr}
\pagestyle{fancy}
\rhead{\small \color{gray} Notițe de Adrian Manea}
\lhead{\small \color{gray} ETTI-AM2, 2017---2018, M. Joița \& A. Niță}


\usepackage[autostyle = false, style = german]{csquotes}
\usepackage[romanian]{babel}
\newcommand\qq{\enquote}

%% COMENZILE MELE
\renewcommand*{\proofname}{Dem.:}
\newcommand\bb{\mathbb}
\newcommand\dr{\mathrm}
\newcommand\kal{\mathcal}
\newcommand\fk{\mathfrak}
\newcommand\op{\oplus}
\newcommand\ot{\otimes}
\newcommand\ds{\displaystyle}
\newcommand\id{\indent}
\newcommand\nid{\noindent}
\newcommand\seq{\subseteq}
\newcommand\sneq{\subsetneq}
\newcommand\speq{\supseteq}
\newcommand\spneq{\supsetneq}
\newcommand\rar{\rightarrow}
\newcommand\Rar{\Rightarrow}
\newcommand\xrar{\xrightarrow}
\newcommand\lar{\leftarrow}
\newcommand\Lar{\Leftarrow}
\newcommand\xlar{\xleftarrow}
\newcommand\lrar{\leftrightarrow}
\newcommand\Lrar{\Leftrightarrow}
\newcommand\ideal{\trianglelefteq}
\newcommand\ai{\text{ a.\^{i}. }}
\newcommand\sk{(\textit{Schi\c{t}\u{a}}) }
\newcommand\rhup{\rightharpoonup}
\newcommand\rhdn{\rightharpoondown}
\newcommand\lhup{\leftharpoonup}
\newcommand\lhdn{\leftharpoondown}
\newcommand\vid{\emptyset}
\newcommand\wh{\widehat}
\newcommand\ol{\overline}
\newcommand\NN{\mathbb{N}}
\newcommand\RR{\mathbb{R}}
\newcommand\ZZ{\mathbb{Z}}
\newcommand\QQ{\mathbb{Q}}
\newcommand\CC{\mathbb{C}}
\newcommand\Rez{\mathrm{Rez}}

%% STILURILE MELE
\newtheoremstyle{dotless}{}{}{\itshape}{}{\bfseries}{:}{ }{}

\theoremstyle{dotless}
\newtheorem{theorem}{Teorem\u{a}}
\newtheorem{proposition}{Propozi\c{t}ie}
\newtheorem{lemma}{Lem\u{a}}[section]
\newtheorem{corollary}{Corolar}
\newtheorem{exercise}{Exerci\c{t}iu}

\newtheoremstyle{dotlessdr}{}{}{}{}{\bfseries}{:}{ }{}
\theoremstyle{dotlessdr}
\newtheorem{example}{Exemplu}
\newtheorem{definition}{Defini\c{t}ie}
\newtheorem{remark}{Observa\c{t}ie}




\begin{document}
% \definecolor{labelkey}{RGB}{222,0,0}

\begin{center}
  {\large\textbf{Seminar 9bis \\
      Teorema reziduurilor (supliment)}}
\end{center}

\vspace{2cm}

\indent\indent Principalele metode de a calcula reziduurile unei funcții
sînt redate în teorema următoare.

\begin{theorem}[Calculul reziduurilor]\label{thm:calcul-rez}
  \begin{enumerate}[(1)]
  \item $ \Rez(f, a) = c_{-1} $, unde $ c_{-1} $ este coeficientului lui $ \dfrac{1}{z - a} $
    din dezvoltarea în serie Laurent a funcției $ f $ în vecinătatea singularității $ z = a $;
  \item Dacă $ z = a $ este pol de ordinul $ p \geq 2 $ pentru $ f $, atunci:
    \[
      \Rez(f, a) = \frac{1}{(p - 1)!} \lim_{z \to a} \Big[ (z - a)^p f(z) \Big]^{(p-1)};
    \]
  \item Dacă $ z = a $ este pol simplu pentru $ f $, atunci, particularizînd formula
    de mai sus, avem:
    \[
      \Rez(f, a) = \lim_{z \to a} (z - a) f(z);
    \]
  \item Dacă $ f $ se poate scrie ca un cît de două funcții $ A, B $, olomorfe
    în jurul punctului $ a $ și dacă $ a $ este pol simplu pentru $ f $, adică
    $ B(a) = 0 $, atunci:
    \[
      \Rez(f, a) = \lim_{z \to a} \frac{A(z)}{B'(z)}.
    \]
  \end{enumerate}
\end{theorem}

\vspace{1cm}

1. Să se calculeze următoarele integrale:
\begin{enumerate}[(a)]
\item $ I = \ds\int_{|z| = 2} \frac{dz}{z^2 - 1} $;
\item $ I = \ds\int_\gamma \frac{dz}{z^4 + 1}, \gamma: x^2 + y^2 - 2x = 0 $;
\item $ I = \ds\int_{|z| = 3} \frac{z^2 + 1}{(z - 1)^2 (z + 2)} dz $.
\end{enumerate}

\emph{Soluție:} (a) Punctele $ z = \pm 1 $ sînt poli de ordinul 1 pentru
funcția $ f(z) = \dfrac{1}{z^2 - 1} $. Ele sînt situate în interiorul discului pe
care integrăm, cu $ |z| = 2 $, deci putem aplica teorema reziduurilor:
\[
  I = 2 \pi i \cdot \Big( \mathrm{Rez}(f, z_1) + \mathrm{Rez}(f, z_2) \Big).
\]

Calculăm separat reziduurile:
\begin{align*}
  \Rez(f, z_1) &= \lim_{z \to 1} (z - 1) \cdot \frac{1}{z^2 - 1} = \frac{1}{2} \\
  \Rez(f, z_2) &= \lim_{z \to -1} (z + 1) \cdot \frac{1}{z^2 - 1} = -\frac{1}{2}.
\end{align*}

Rezultă:
\[
  I = \int_{|z| = 2} \frac{dz}{z^2 - 1} = 2 \pi i \cdot \Big( \frac{1}{2} - \frac{1}{2} \Big) = 0.
\]

\vspace{1cm}

(b) Curba $ \gamma $ este un cerc centrat în $ (1, 0) $ și cu raza $ 1 $. Căutăm
polii funcției $ f(z) = \frac{1}{z^4 + 1} $ care se află în interiorul lui $ \gamma $.

Avem succesiv:
\begin{align*}
  z^4 + 1 &= 0 \Rightarrow z^4 = -1 = \cos \pi + i \sin \pi \Rightarrow \\
  z &= \sqrt[4]{\cos \pi + i \sin \pi} \Rightarrow \\
  z_k &= \cos \frac{\pi + 2k\pi}{4} + i \sin \frac{\pi + 2k\pi}{4}, k = 0, 1, 2, 3.
\end{align*}

Doar punctele $ z_0, z_3 $ se află în interiorul discului delimitat de $ \gamma $
și calculăm reziduurile în aceste puncte.

Putem aplica formula din Teorema \ref{thm:calcul-rez} (4) și avem:
\begin{align*}
  \Rez(f, z_0) &= \frac{A(z)}{B'(z)}\Big|_{z_0} = \frac{1}{4z^3}\Big|_{z_0} %
                 = -\frac{1}{4} e^{\frac{\pi i}{4}} \\
  \Rez(f, z_3) &= \frac{A(z)}{B'(z)}\Big|_{z_3} = \frac{1}{4z^3}\Big|_{z_3} %
                 = - \frac{1}{4} e^{\frac{7 \pi i}{4}}.
\end{align*}

Rezultă:
\[
  I = \int_\gamma \frac{dz}{z^4 + 1} = %
  2 \pi i \Big( -\frac{1}{4} e^{\frac{\pi i}{4}} - \frac{1}{4} e^{\frac{7 \pi i}{4}} \Big) = %
  -\frac{\pi \sqrt{2} i}{2}.
\]

\vspace{1cm}

(c) Avem doi poli, $ z = 1, z = -2 $ în interiorul conturului. Se vede că $ z_1 = 1 $
este pol de ordinul 2, iar $ z_2 = -2 $ este pol de ordinul 1. Calculăm reziduurile:
\begin{align*}
  \Rez(f, z_1) &= \frac{1}{2!} \lim_{z \to 1} \Big[ (z - 1)^2 \frac{z^2 + 1}{(z - 1)^2(z + 2)} \Big] \\
               &= \frac{1}{2!} \lim_{z \to 1} \frac{z^2 + 4z - 1}{(z + 2)^2} \\
               &= \frac{2}{9} \\
  \Rez(f, z_2) &= \lim_{z \to -2} (z + 2) \frac{z^2 + 1}{(z - 1)^2(z + 2)} = \frac{5}{9}.
\end{align*}

Rezultă:
\[
  I = \int_{|z| = 3} \frac{z^2 + 1}{(z - 1)^2(z + 2)} dz = \frac{14}{9} \pi i.
\]

\vspace{1cm}

2. Să se calculeze integralele:
\begin{enumerate}[(a)]
\item $ \ds\int_{|z| = 1} \frac{dz}{\sin z} $;
\item $ \ds\int_{|z| = 1} \frac{e^z}{z^2}dz $;
\item $ \ds\int_{|z| = 5} z e^{\frac{3}{z}} dz $;
\item $ \ds\int_{|z - 1| = 1} \frac{dz}{(z + 1)(z - 1)^3} $;
\item $ \ds\int_{|z| = 1} \frac{\sin z}{z^4} dz $.
\end{enumerate}

\emph{Indicații}: (a) $ z = 0 $ este singurul pol din interiorul domeniului;

(b), (c) Dezvoltăm în serie Laurent și identificăm reziduurile folosind
Teorema \ref{thm:calcul-rez} (1).

(d) Avem $ z = 1 $ pol de ordin 3 și $ z = -1 $ pol simplu. Doar $ z = 1 $ se află
în interiorul domeniului și dezvoltăm în serie Laurent după puterile lui $ z - 1 $:
\begin{align*}
  \frac{1}{z + 1} &= \frac{1}{2 - (-(z - 1))} \\
                  &= \frac{1}{2} \frac{1}{1 - \big(-\frac{z - 1}{2}\big)} \\
                  &= \frac{1}{2} \sum_{n \geq 0} (-1)^n \frac{(z - 1)^n}{2^n} \\
                  &= \sum_{n \geq 0} (-1)^n \frac{(z - 1)^n}{2^{n+1}},
\end{align*}
pentru $ |z - 1| < 2 $.

Rezultă:
\[
  \frac{1}{(z + 1)(z - 1)^3} = \sum_{n \geq 0} (-1)^n \frac{(z - 1)^{n-3}}{2^{n+1}} = %
  \frac{1}{2(z - 1)^3} - \frac{1}{4(z - 1)^2} + \frac{1}{8(z - 1)} - \frac{1}{16} + \dots,
\]
deci $ \Rez(f, 1) = \frac{1}{8} $.

(e) Avem $ z = 0 $ pol de ordinul 3. Putem calcula reziduul folosind dezvoltarea în
serie Laurent sau cu formula din Teorema \ref{thm:calcul-rez}(2).

\vspace{2cm}

\textbf{Aplicații ale teoremei reziduurilor}

Putem folosi teorema reziduurilor pentru a calcula integrale trigonometrice de forma:
\[
  I = \int_0^{2 \pi} R(\cos \theta, \sin \theta) d \theta,
\]
unde $ R $ este o funcție rațională.

Facem schimbarea de variabilă $ z = e^{i\theta} $ și atunci, pentru $ \theta \in [0, 2\pi] $,
$ z $ descrie cercul $ |z| = 1 $, o dată, în sens direct.

Folosim formulele lui Euler:
\begin{align*}
  \cos \theta &= \frac{e^{i\theta} + e^{-i\theta}}{2} = \frac{1}{2} \Big( z + \frac{1}{z} \Big) \\
  \sin \theta &= \frac{e^{i\theta} - e^{-i\theta}}{2i} = \frac{1}{2i} \Big( z - \frac{1}{z} \Big).
\end{align*}

Atunci, dacă $ z = e^{i\theta} $, rezultă $ dz = ie^{i\theta} d\theta = izd\theta $,
iar integrala devine:
\[
  I = \int_0^{2 \pi} R(\cos \theta, \sin \theta) d\theta = \oint_{|z| = 1} R_1(z) dz,
\]
unde:
\[
  R_1(z) = \frac{1}{iz} R\Big( \frac{z^2 + 1}{2z}, \frac{z^2 - 1}{2iz} \Big).
\]

Această funcție poate avea poli și deci putem folosi teorema reziduurilor. Dacă
$ a_1, \dots, a_n $ sînt polii din interiorul cercului unitate, avem:
\[
  I_1 = 2 \pi i \sum_{k \geq 1} \Rez(R_1, a_k).
\]

\vspace{1cm}

Să vedem cîteva exemple:
\begin{enumerate}[(a)]
\item $ \ds\int_0^{2 \pi} \frac{d\theta}{2 + \cos \theta} $;
\item $ \ds\int_0^{2 \pi} \frac{d\theta}{1 + 3 \cos^2 \theta} $;
\item $ \ds\int_0^{2 \pi} \frac{\sin \theta}{1 + i \sin \theta} d\theta $;
\item $ \ds\int_0^{2 \pi} \frac{d\theta}{1 + a \sin \theta}, |a| < 1, a \in \RR $.
\end{enumerate}

\emph{Soluție:}

(a) Notăm $ z = e^{i\theta} $, cu $ \theta \in [0, 2\pi] $. Atunci avem succesiv:
\begin{align*}
  dz &= ie^{i\theta} d\theta = izdz \Rightarrow d\theta = \frac{dz}{iz} \\
  \cos \theta &= \frac{e^{i\theta} + e^{-i\theta}}{2} = \frac{1}{2} \Big( z + \frac{1}{z} \Big) \\
  \frac{1}{2 + \cos \theta} &= \frac{2z}{z^2 + 4z + 1} \\
  \int_0^{2 \pi} \frac{d\theta}{2 + \cos \theta} &= \oint_{|z| = 1} \frac{2z}{z^2 + 4z + 1} \frac{dz}{iz} \\
     &= -2i \oint_{|z| = 1} \frac{dz}{z^2 + 4z + 1}.
\end{align*}

Acum folosim teorema reziduurilor. Singularitățile funcției $ f(z) = \frac{1}{z^2 + 4z + 1} $
sînt $ z = -2 \pm \sqrt{3} $, care sînt poli simpli. Numai $ z = -2 + \sqrt{3} $ se află
în interiorul cercului $ |z| = 1 $ și calculăm reziduul folosind
Teorema \ref{thm:calcul-rez}(2).

\end{document}
